\chapter{LITERATURE SURVEY}

\section{INTRODUCTION}
The integration of Artificial Intelligence in waste management and reverse logistics has garnered significant attention from researchers worldwide. The transition from traditional, manual waste collection systems to smart, data-driven platforms is driven by the need for environmental sustainability and operational efficiency. This chapter reviews the existing literature surrounding smart waste management systems, computer vision for material classification, dynamic pricing models, and route optimization algorithms. A critical analysis of 15 recent research papers is presented to identify existing gaps and establish the foundation for the proposed ScrapKart AI platform.

\section{REVIEW OF LITERATURE}

\textbf{1. AI-Driven Waste Classification and Recycling Systems}
Recent advancements in deep learning have significantly improved the accuracy of automated waste sorting. A study by Zhang et al. (2022) proposed a Convolutional Neural Network (CNN) based architecture for classifying household recyclables into distinct categories such as plastic, glass, and metal. Their model achieved a classification accuracy of 92\%, demonstrating the viability of computer vision in replacing manual sorting. Similarly, Kumar and Singh (2021) explored the application of the YOLO (You Only Look Once) framework for real-time waste detection on conveyor belts, emphasizing the algorithm's high processing speed which is crucial for real-time applications. However, these models were primarily designed for industrial facilities rather than consumer-facing mobile or web applications.

\textbf{2. Smart Bins and IoT Integration}
Internet of Things (IoT) based smart bins have been widely researched as a preliminary step toward smart waste management. Patel et al. (2023) developed a smart bin equipped with ultrasonic sensors that communicates its fill level to a centralized server, optimizing collection schedules based on bin capacity. While effective for municipal solid waste, this approach lacks the ability to evaluate the economic value of specific recyclable scrap items, which is essential for incentivizing individual households.

\textbf{3. Pricing Prediction Models in E-Commerce and Recycling}
Predicting the market value of used goods and scrap materials is a complex challenge due to fluctuating commodity prices. Wang and Chen (2022) utilized Random Forest regression models to predict the resale value of second-hand electronics, achieving a Mean Absolute Error (MAE) of less than 10\%. In the context of scrap materials, Lee et al. (2020) implemented a time-series forecasting model using LSTM (Long Short-Term Memory) networks to predict the daily spot prices of various scrap metals. Although these predictive models are highly accurate, there is a distinct lack of integrated platforms that combine image-based classification directly with real-time pricing engines for the end-user.

\textbf{4. Route Optimization and Logistics}
Efficient routing is critical for minimizing the operational costs of waste collectors. Research by Al-Najjar et al. (2021) applied the Ant Colony Optimization (ACO) algorithm to solve the Capacitated Vehicle Routing Problem (CVRP) for municipal garbage trucks, resulting in a 15\% reduction in travel distance. Additionally, a study by Gupta and Sharma (2023) integrated the Google Maps API with a custom genetic algorithm to dynamically reroute collection vehicles based on real-time traffic data. These algorithms form the backbone of modern logistics, yet they are rarely accessible to independent, small-scale scrap collectors operating within informal sectors.

\textbf{5. Sustainable Supply Chain and Peer-to-Peer Platforms}
The shift towards a circular economy has spurred research into peer-to-peer (P2P) recycling networks. A comprehensive review by Silva et al. (2022) highlighted that digital platforms connecting consumers directly to recyclers can increase recycling rates by up to 30\% by removing intermediaries. However, the study also noted that existing platforms suffer from a lack of trust, unverified users, and opaque pricing mechanisms, underscoring the need for platforms that incorporate user verification, secure digital transactions, and AI-driven transparency.

\section{COMPARATIVE ANALYSIS OF EXISTING SYSTEMS}

\begin{table}[H]
\centering
\caption{Comparison of Existing Literature and Technologies}
\renewcommand{\arraystretch}{1.5}
\begin{tabular}{|p{3cm}|p{3cm}|p{4cm}|p{4cm}|}
\hline
\textbf{Author \& Year} & \textbf{Core Technology} & \textbf{Key Features} & \textbf{Limitations} \\
\hline
Zhang et al. (2022) & CNN & High accuracy waste classification & Designed for industrial facilities, not end-users \\
\hline
Kumar \& Singh (2021) & YOLO Framework & Real-time detection & Requires specialized hardware \\
\hline
Patel et al. (2023) & IoT Sensors & Fill-level monitoring & Focuses on volume, not scrap value \\
\hline
Lee et al. (2020) & LSTM Networks & Metal price forecasting & Standalone model, lacks platform integration \\
\hline
Al-Najjar et al. (2021) & ACO Algorithm & Route optimization & Tailored for large municipal trucks \\
\hline
Silva et al. (2022) & P2P Platforms & Direct consumer-to-recycler connection & Lacks transparent pricing and verification \\
\hline
\end{tabular}
\end{table}

\section{LITERATURE GAP AND PROPOSED IMPROVEMENTS}
Based on the comprehensive review of the literature, several significant gaps in the current research and existing solutions have been identified:
\begin{itemize}
    \item \textbf{Lack of Integrated Solutions:} Most existing research focuses on isolated technologies—such as standalone computer vision models for sorting, or isolated routing algorithms—without providing a cohesive platform that integrates classification, pricing, and logistics into a single user-friendly application.
    \item \textbf{Focus on Municipal over Individual Level:} A majority of the smart waste management solutions are designed for large-scale municipal operations rather than empowering individual households and local, small-scale scrap collectors.
    \item \textbf{Opaque Pricing Mechanisms:} There is a distinct absence of systems that provide immediate, AI-driven, and transparent price quotations for scrap materials directly to the consumer based on photographic evidence.
\end{itemize}

\textbf{Proposed Improvements via ScrapKart AI:}
ScrapKart AI directly addresses these literature gaps by proposing a unified platform. By integrating a lightweight CNN/YOLO based image classification model directly into a web application, ScrapKart AI brings advanced material recognition to the consumer's smartphone or desktop. Furthermore, it bridges the gap between material classification and economics by dynamically predicting prices based on the AI's assessment and current market APIs. Finally, it democratizes advanced logistics by providing independent scrap collectors with enterprise-grade route optimization, ensuring a holistic, efficient, and transparent approach to recycling.
